\documentclass[conference]{IEEEtran}
\IEEEoverridecommandlockouts

\usepackage{amsmath,amssymb,amsfonts}
\usepackage{graphicx}
\usepackage{booktabs}
\usepackage{multirow}
\usepackage{cite}
\usepackage{algorithm}
\usepackage{algorithmic}
\usepackage{url}
\usepackage{xcolor}
\usepackage{float}
\usepackage[hidelinks]{hyperref}

\floatstyle{ruled}
\restylefloat{algorithm}

\graphicspath{{../figure/}}

\newcommand{\ordi}{ORDI}
\newcommand{\eg}{e.g.,}
\newcommand{\ie}{i.e.,}
\renewcommand{\IEEEkeywordsname}{Keywords}

\begin{document}

\title{\ordi{}: Orbit-Aware Redundant Distributed\\Inference for LEO Earth Observation Constellations}

\author{
\IEEEauthorblockN{An-Che Liang, Guan-Ming Chiu, Pei-Yu Wu}
\IEEEauthorblockA{National Taiwan University}
}

\maketitle

\begin{abstract}
\end{abstract}

\begin{IEEEkeywords}
LEO constellations, Earth observation, distributed inference, fault-tolerant scheduling
\end{IEEEkeywords}

\setcounter{section}{1}
\section{Background and Related Work}
\label{sec:motivation}

Orbital edge computing treats a satellite constellation as a distributed
computing platform rather than a collection of sensors that merely forward raw
data. Denby and Lucia establish this design space and expose the coupled
constraints of orbital motion, onboard compute, energy, sensing, and
downlink~\cite{denby2020oec}. Subsequent systems provide complementary
mechanisms: Krios places and migrates persistent services across moving LEO
resources~\cite{krios2024}, LEOEdge coordinates heterogeneous satellite and
ground resources for AI inference~\cite{leoedge2025}, and SECO jointly schedules
wide-area observation, routing, and compute placement~\cite{seco}. These systems
primarily target service availability, latency, coverage, or resource
efficiency under nominal mobility.

Deadline-aware mobile-edge schedulers provide useful methods for dynamic
arrivals and constrained communication. Liu et al.\ minimize the deadline
violation ratio for dependent tasks through scheduling, migration, and
merging~\cite{liu2023deadline}, while TODG makes distributed online offloading
decisions under hard delay constraints and stochastic communication
resources~\cite{yue2022todg}. These formulations motivate deadline misses,
rather than average latency alone, as a reliability metric.

Selective replication is a separate approach to reducing delay and execution
risk. Learning-based vehicular-edge schemes estimate worker quality and
replicate tasks across a chosen set of workers~\cite{sun2018replication}; their
analysis further shows that excessive replication can increase queue contention
and perform worse than a smaller replica count~\cite{sun2021replication}.

\ordi{} combines these three lines of work in an online scheduler for LEO
Earth-observation inference. It splits requests into spatial tasks, places them
on queue-aware helpers, and selects contact-aware routes that can complete
computation and ground delivery before an absolute deadline. It also learns
failure-domain risk from execution outcomes and selectively adds fault-disjoint
backups when their expected reliability gain justifies the additional compute
and communication cost. This integration makes deadline feasibility, orbital
resource variation, and correlated failures part of one scheduling decision.

\section{Modeling}
\label{sec:modeling}

\subsection{System and Task Model}
\label{sec:model}

Orbital contacts, compute queues, and energy--thermal states evolve over time,
whereas continuously recomputing the schedule after every state change would
incur prohibitive control and optimization overhead. We therefore adopt a
rolling-horizon discretization in which scheduling decisions are revised at
fixed epochs indexed by $\ell\in\mathbb{Z}_{\ge 0}$. Each epoch has duration
$\Delta$ and interval
\begin{equation}
 \mathcal{I}_{\ell}=[t_{\ell},t_{\ell+1}),
 \qquad t_{\ell}=\ell\Delta .
 \label{eq:epoch}
\end{equation}
At $t_{\ell}$, the scheduler refreshes the satellite states and pending task
set, and constructs a rolling-horizon plan using contacts predicted through
the relevant task deadlines. The epoch length determines the replanning
cadence; it does not restrict an operation to the current interval, and
resource reservations may extend beyond $\mathcal{I}_{\ell}$. Hereafter, $t$
denotes the current decision time $t_{\ell}$.

The constellation is represented by a time-varying directed contact graph
$\mathcal{G}(t)=(\mathcal{N},\mathcal{E}(t))$, where
$\mathcal{N}=\mathcal{S}\cup\mathcal{D}$ contains satellites
$\mathcal{S}$ and ground stations $\mathcal{D}$. A contact
$e=(i,j,o_e,c_e,R_e,\rho_e,\kappa_e)$ denotes a directed communication
opportunity from $i$ to $j$, where $o_e$ and $c_e$ are its opening and closing
times, respectively. During the interval $[o_e,c_e]$, the contact has rate
$R_e$, success probability $\rho_e$, and type
$\kappa_e\in\{\mathrm{ISL},\mathrm{downlink}\}$. Its raw capacity is
$K_e=R_e(c_e-o_e)$ bits. A feasible store-and-forward route is a
time-ordered sequence of contacts and may include waiting at intermediate
satellites.

At decision time $t$, satellite $i$ is characterized by
\begin{equation}
  \boldsymbol{\sigma}_i(t)
  =\bigl(C_i(t),B_i(t),\Theta_i(t),Q_i(t),A_i(t),\rho_i(t)\bigr),
  \label{eq:satellite-state}
\end{equation}
where $C_i$ is its effective compute rate in FLOP/s, $B_i$ is battery energy,
$\Theta_i$ is payload temperature, $Q_i$ is already-queued work in FLOPs,
$A_i\in\{0,1\}$ is availability, and $\rho_i$ is node reliability. The
effective rate $C_i$ incorporates thermal throttling. Let $H_i(t)$ indicate
whether the platform is operational. Availability is defined as
\begin{equation}
 A_i(t)=
 \mathbb{1}\!\left[
   H_i(t)=1,\ B_i(t)\ge B_i^{\min},\
   \Theta_i(t)<\Theta_i^{\max}
 \right].
 \label{eq:availability}
\end{equation}
The battery and thermal dynamics update $C_i$ and $A_i$ between consecutive
decision epochs.

Each EO request $k$ is represented as
$\mathcal{T}_k=(s_k,r_k,D_k,\mathcal{V}_k)$, where $s_k$ is the imaging
satellite, $r_k$ is the release time, $D_k$ is the absolute delivery deadline,
and $\mathcal{V}_k$ is a grid of image tiles. Each tile $v\in\mathcal{V}_k$ is
described by
\begin{equation}
  \tau_{kv}=(d^{\mathrm{in}}_{kv},d^{\mathrm{out}}_{kv},w_{kv}),
  \label{eq:tile}
\end{equation}
where $d^{\mathrm{in}}_{kv}$ and $d^{\mathrm{out}}_{kv}$ are input and result
sizes in bits, and $w_{kv}$ is inference work in FLOPs.
The two data volumes are distinguished because onboard inference transforms a
high-dimensional image tile into a task-dependent product, such as a class
label, bounding boxes, or a segmentation mask.

Orbital motion is predictable; therefore each satellite is assumed to know
the contact graph over the scheduling horizon. This knowledge specifies when
links will exist, but not the contemporaneous resource state of another
satellite. At time $t$, satellite $i$ directly observes only
$\boldsymbol{\sigma}_i(t)$. For $j\ne i$, its most recent estimate is
\begin{equation}
 \widehat{\boldsymbol{\sigma}}^{\,i}_j(t)
 =\boldsymbol{\sigma}_j(t^i_j),
 \qquad
 a^i_j(t)=t-t^i_j,
 \label{eq:state-estimate}
\end{equation}
where $t^i_j$ is the generation time of the latest state advertisement
received by $i$, and $a^i_j$ is its age. The state of $j$ is unknown to $i$
until such an advertisement is received.

Neighboring satellites exchange state advertisements and task data over
inter-satellite links. To limit flooding overhead and preserve contact
capacity for task traffic, each epoch's state advertisement is forwarded over
at most three ISL hops. Let $\mathcal{A}^{(h)}_i(e)$ be the advertisements
available to satellite $i$ after at most $h$ forwarding hops in epoch $e$, and
let $\mathcal{N}_i(e)$ be its incoming ISL neighbors during that epoch. The
exchange can be summarized as
\begin{align}
 \mathcal{A}^{(0)}_i(e)
   &=\{(i,\boldsymbol{\sigma}_i(t_e),t_e)\}, \nonumber\\
 \mathcal{A}^{(h)}_i(e)
   &=\operatorname{fresh}\!\left(
       \mathcal{A}^{(h-1)}_i(e)
       \cup
       \bigcup_{j\in\mathcal{N}_i(e)}
       \mathcal{A}^{(h-1)}_j(e)
     \right), \nonumber\\[-2pt]
 &\hspace{24mm} h=1,2,3.
 \label{eq:advertisement-exchange}
\end{align}
where $\operatorname{fresh}(\cdot)$ retains only the newest advertisement per
origin. A received entry becomes usable only after its contact-constrained
arrival time. This bounded dissemination extends state awareness beyond
immediate neighbors without requiring network-wide flooding. Task data may
traverse multiple ISLs, and a downlink delivers the resulting data to a ground
station.
Control and task traffic share the finite capacities of the corresponding
contact windows. Thus, knowledge of future connectivity is global and
deterministic, whereas
knowledge of remote satellite state is local, delayed, and acquired through
communication.

\subsection{Constraints and Problem Formulation}
\label{sec:formulation}

For each tile $v\in\mathcal{V}_k$, the decisions include a split width
$q_{kv}\in\{1,2,4\}$ and a reconstruction-group count
$n_{kv}\in\{1,2\}$. Following similar distributed-inference work
\cite{zhao2018deepthings,tang2022partitioning}, we allow at most four shards
because over-sharding increases communication, synchronization, and scheduling
overhead.

A reconstruction group is an end-to-end plan that independently completes tile
$v$, indexed by $g\in\{1,\ldots,n_{kv}\}$. We choose $n_{kv}\le2$ because
each additional group schedules another $q_{kv}$ shard placements, consuming
compute and link capacity while enlarging the combinatorial placement search.

For group $g$, shard $m\in\{1,\ldots,q_{kv}\}$ carries
$d^{\mathrm{in}}_{kv}/q_{kv}$ input bits, requires $w_{kv}/q_{kv}$ FLOPs, and
produces $d^{\mathrm{out}}_{kv}/q_{kv}$ result bits. Its placement
\begin{equation}
 x_{kvgm}=(h_{kvgm},P^{\mathrm{in}}_{kvgm},
 P^{\mathrm{gnd}}_{kvgm})
\label{eq:placement}
\end{equation}
where $h_{kvgm}$ is the compute satellite,
$P^{\mathrm{in}}_{kvgm}$ and $P^{\mathrm{gnd}}_{kvgm}$ are the
time-respecting input and output paths, respectively. Each path contains at
most $H_{\max}^{\mathrm{ISL}}=12$ ISL hops.

\noindent Let $F^{\mathrm{in}}_{kvgm}$ be the time at which the input shard reaches its
compute satellite, and let $S^{\mathrm{cmp}}_{kvgm}$ be its computation start times.
The shard computation is constrained by
\begin{equation}
 \begin{aligned}
 S^{\mathrm{cmp}}_{kvgm}
 &\ge F^{\mathrm{in}}_{kvgm},\\
 \int_t^{S^{\mathrm{cmp}}_{kvgm}}
 C_{h_{kvgm}}(\tau)\,d\tau
 &\ge Q_{h_{kvgm}}(t),\\
 \int_{S^{\mathrm{cmp}}_{kvgm}}^{F^{\mathrm{cmp}}_{kvgm}}
 C_{h_{kvgm}}(\tau)\,d\tau
 &=\frac{w_{kv}}{q_{kv}}.
 \end{aligned}
 \label{eq:compute-time}
\end{equation}
The result path starts after computation. Let $F_{kvgm}$ denote the time at
which the result reaches the ground. A group completes after all $q_{kv}$ of its
results arrive, while a tile completes when any selected reconstruction group
completes:
\begin{equation}
 T_{kvg}=\max_{1\le m\le q_{kv}}F_{kvgm},
 \qquad
 T_{kv}=\min_{1\le g\le n_{kv}}T_{kvg}.
 \label{eq:completion-time}
\end{equation}
Every admitted placement obeys
\begin{equation}
 \begin{aligned}
 r_k&\le t,\qquad F_{kvgm}\le D_k,\\
 A_i(t)&=1,\quad
 \forall i\in
 \mathcal{S}(P^{\mathrm{in}}_{kvgm}\cup
 P^{\mathrm{gnd}}_{kvgm}).
 \end{aligned}
 \label{eq:deadline-availability}
\end{equation}
Here, $\mathcal{S}(P)$ denotes the satellite nodes traversed by path $P$.
At decision time, a satellite evaluates remote availability using its latest
estimate from Equation~\eqref{eq:state-estimate}; stale information can
therefore cause a planned placement to fail during execution.

Shared resources are explicitly serialized. For every contact window $e$,
\begin{equation}
 \sum_{x_{kvgm}\,\ni\,e} b_{kvgm,e}\le K_e,
 \label{eq:contact-capacity}
\end{equation}
and each transmission interval must lie inside $[o_e,c_e]$. Intervals using
the same satellite terminal cannot overlap, and compute intervals on one
compute satellite cannot overlap. State advertisements and task transfers
consume capacity from the same contact-window budget.

Within a group, compute satellites are distinct. A backup group must
additionally avoid the primary compute satellites, their orbital planes, and
satellite nodes already used by the primary routes. If $\delta(i)$ denotes the
fault domain of satellite $i$, then for primary group $1$ and each selected
backup group $2\le g\le n_{kv}$,
\begin{equation}
\delta(h_{kv1m})\ne\delta(h_{kvgm'})
\quad \forall\,m,m'\in\{1,\ldots,q_{kv}\},
\label{eq:fault-disjoint}
\end{equation}
with the analogous node-disjoint condition on their routes.

Let $\mathcal{C}_{kvg}$ be the physical node and link components required by
group $g$. Its success probability is
$p_{kvg}=\prod_{c\in\mathcal{C}_{kvg}}\rho_c$. Because groups can share a
source or contact component, tile reliability uses inclusion--exclusion rather
than assuming independent replicas:
\begin{equation}
 P_{kv}=
 \sum_{\emptyset\ne H\subseteq\{1,\ldots,n_{kv}\}}
 (-1)^{|H|+1}
 \prod_{c\in\bigcup_{g\in H}\mathcal{C}_{kvg}}\rho_c ,
 \label{eq:tile-reliability}
\end{equation}
Let $L_{kv}=T_{kv}-t$ be the service latency, and let $Z_{kv}$ be the total
routed bits counted per hop. The online scheduling problem is
\begin{equation}
 \max_{\{q_{kv},n_{kv},x\}}
 \sum_{k}\sum_{v\in\mathcal{V}_k}
 \left[
 P_{kv}e^{-\alpha L_{kv}}
 -\lambda_C Z_{kv}
 -\lambda_R(n_{kv}-1)
 \right]
\label{eq:objective}
\end{equation}
subject to Equations~\eqref{eq:compute-time}--\eqref{eq:fault-disjoint}, with
$q_{kv}\in\{1,2,4\}$ and $n_{kv}\in\{1,2\}$. The physical model separately
integrates compute and radio energy and updates $B_i$ and $\Theta_i$; therefore
energy affects subsequent feasibility through Equation~\eqref{eq:availability}
rather than through an optimistic per-placement estimate.

\bibliographystyle{IEEEtran}
\bibliography{references}

\end{document}
